\documentclass[a4paper,11pt]{article}
\usepackage[utf8]{inputenc}
\usepackage[bahasa]{babel}
\usepackage{geometry}
\usepackage{graphicx}
\usepackage{array}
\usepackage{multirow}
\usepackage{longtable}
\usepackage{booktabs}
\usepackage{enumitem}

% Pengaturan Margin
\geometry{
    top=2.5cm,
    bottom=2.5cm,
    left=3cm,
    right=2.5cm
}

% Definisi kolom tabel
\newcolumntype{L}[1]{>{\raggedright\let\newline\\\arraybackslash\hspace{0pt}}m{#1}}
\newcolumntype{C}[1]{>{\centering\let\newline\\\arraybackslash\hspace{0pt}}m{#1}}

\begin{document}

% --- KOP SURAT ---
\begin{center}
    \textbf{\large INSTITUT TEKNOLOGI BANDUNG}\\
    \textbf{\large SEKOLAH TEKNIK ELEKTRO DAN INFORMATIKA}\\
    \textbf{PROGRAM STUDI MAGISTER INFORMATIKA}\\
    \vspace{0.2cm}
    \hrule height 1pt
    \vspace{0.1cm}
    \hrule height 0.5pt
    \vspace{0.5cm}
    \textbf{\large LEMBAR EVALUASI DIRI MAHASISWA}\\
    \textbf{Semester I - 2025/2026} % Disesuaikan dengan periode saat ini
\end{center}
\vspace{0.5cm}

% --- DATA MAHASISWA ---
\section*{A. Data Mahasiswa}
\begin{table}[h!]
    \renewcommand{\arraystretch}{1.3}
    \begin{tabular}{L{4cm} L{0.5cm} L{10cm}}
        \textbf{Nama} & : & Jatniko Nur Mutaqin \\
        \textbf{NIM} & : & 23523314 \\
        \textbf{Jalur Pilihan} & : & Smart System (Smart-X) \\
        \textbf{Pembimbing} & : & Ir. I Gusti Bagus Baskara Nugraha, ST, MT, Ph.D \\
        \textbf{Judul Tesis} & : & Restorasi Dokumen Terdegradasi Menggunakan Generative Adversarial Network dengan Diskriminator Dual-Modal dan Optimasi Loss Function Berorientasi HTR
    \end{tabular}
\end{table}

% --- ABSTRAK ---
\section*{B. Ringkasan Tesis}
\noindent \textbf{Abstrak:} \\
Penelitian ini bertujuan mengembangkan kerangka kerja restorasi dokumen berbasis \textit{Generative Adversarial Network} (GAN) yang secara eksplisit mengintegrasikan umpan balik semantik guna meningkatkan kinerja \textit{Handwritten Text Recognition} (HTR). Kebaruan penelitian terletak pada penerapan strategi \textit{frozen recognizer} yang dikombinasikan dengan fungsi kerugian multiobjektif lima komponen, serta evaluasi kritis terhadap arsitektur diskriminator dual-modal. Metode melibatkan pelatihan generator U-Net yang dimodifikasi, yang belajar dari perbedaan piksel citra dan umpan balik gradien model pengenal teks beku (\textit{frozen}). Hasil evaluasi menggunakan 712 citra uji semisintetis menunjukkan model mampu menurunkan \textit{Character Error Rate} (CER) signifikan dari 83,4\% menjadi 34,9\% ($p<0,001$), mendekati batas bawah teoretis (34,1\%). Model juga mempertahankan kualitas visual tinggi dengan PSNR 30,74 dB dan SSIM 0,987. Studi ini membuktikan bahwa stabilitas sinyal pembelajaran melalui \textit{frozen recognizer} lebih vital daripada kompleksitas arsitektur diskriminator dalam restorasi dokumen historis.

% --- LAPORAN MINGGUAN ---
\section*{C. Laporan Kemajuan Mingguan}
Berikut adalah rekapitulasi kegiatan penelitian dan bimbingan yang telah dilakukan:

\begin{longtable}{|C{1cm}|C{3cm}|L{10.5cm}|}
\hline
\textbf{No} & \textbf{Tanggal} & \textbf{Kegiatan / Kemajuan / Substansi Bimbingan} \\
\hline
\endfirsthead
\hline
\textbf{No} & \textbf{Tanggal} & \textbf{Kegiatan / Kemajuan / Substansi Bimbingan} \\
\hline
\endhead

1 & 12 Sep 2025 & \textbf{Penyusunan Bab I Pendahuluan:}
\begin{itemize}[leftmargin=*, noitemsep]
    \item Diskusi mengenai urgensi metode GAN dalam latar belakang.
    \item Revisi judul menjadi deskriptif (menghapus singkatan di awal) dan perbaikan alur logika pemilihan metode penelitian.
    \item Penyesuaian format penulisan (pengurangan penggunaan bold).
\end{itemize} \\
\hline

2 & 26 Sep 2025 & \textbf{Standardisasi Format Dokumen:}
\begin{itemize}[leftmargin=*, noitemsep]
    \item Konversi format \textit{bullet points} menjadi narasi/penomoran sesuai standar tesis.
    \item Proofreading kata sambung ("Sedangkan") dan konsistensi font Daftar Isi.
    \item Koreksi nomenklatur Program Studi "Magister Informatika".
\end{itemize} \\
\hline

3 & 17 Okt 2025 & \textbf{Penulisan Bab III \& IV (Metodologi \& Perancangan):}
\begin{itemize}[leftmargin=*, noitemsep]
    \item Finalisasi persamaan matematis (rumus fungsi loss dan arsitektur GAN).
    \item Perbaikan notasi matematika (simbol perkalian) dan penomoran bab pada rumus (misal: III.2).
    \item Melengkapi sitasi dan \textit{caption} pada gambar arsitektur sistem.
\end{itemize} \\
\hline

4 & 24 Okt 2025 & \textbf{Strategi Publikasi Ilmiah:}
\begin{itemize}[leftmargin=*, noitemsep]
    \item Analisis target jurnal: Diskusi opsi \textit{IEEE Access} vs \textit{Wiley}.
    \item Pengecekan syarat APC dan kuota klaim institusi.
    \item Ekspansi daftar pustaka untuk memenuhi syarat minimal 40 referensi artikel jurnal bereputasi.
\end{itemize} \\
\hline

5 & 31 Okt 2025 & \textbf{Finalisasi Metadata Paper:}
\begin{itemize}[leftmargin=*, noitemsep]
    \item Koreksi afiliasi penulis ("Graduate Program in Informatics") dan jabatan pembimbing.
    \item Penambahan detail \textit{expertise} (Smart City, Blockchain) pada biodata penulis.
    \item Pengaturan peran \textit{corresponding author} untuk submisi.
\end{itemize} \\
\hline

6 & 14 Nov 2025 & \textbf{Revisi Bab II Tinjauan Pustaka \& Sitasi:}
\begin{itemize}[leftmargin=*, noitemsep]
    \item Eliminasi paragraf duplikat pada naskah paper/tesis.
    \item Revisi istilah "Core Contributions" menjadi "Main Contributions".
    \item Konversi manajemen referensi manual ke BibTeX untuk memperbaiki sitasi yang hilang (Ref 3-6, 20-21) dan duplikasi rumus (Eq 17).
\end{itemize} \\
\hline

7 & 17 Nov 2025 & \textbf{Submisi Makalah:}
\begin{itemize}[leftmargin=*, noitemsep]
    \item Mengirimkan makalah ke jurnal \textit{Advances in Human-Computer Interaction} (Wiley).
    \item Status: \textit{Submitted}.
\end{itemize} \\
\hline

8 & 24 Nov 2025 & \textbf{Penyelesaian Draf Tesis V1:}
\begin{itemize}[leftmargin=*, noitemsep]
    \item Penyerahan Draf Lengkap Tesis (Bab I - VI) hasil penelitian.
    \item Pengurusan administrasi formulir izin usulan seminar tesis.
    \item Persiapan materi slide presentasi untuk seminar hasil.
\end{itemize} \\
\hline

9 & 28 Nov 2025 & \textbf{Revisi Akhir \& Helicopter View:}
\begin{itemize}[leftmargin=*, noitemsep]
    \item Memperjelas benang merah penelitian (\textit{helicopter view}) dari masalah ke solusi.
    \item Konversi narasi metodologi yang kompleks menjadi tabel agar sistematis.
    \item Detail teknis spesifikasi citra (dimensi/resolusi) dan perbaikan format (italic istilah asing, desimal).
\end{itemize} \\
\hline

\end{longtable}

% --- LUARAN ---
\section*{D. Status Luaran Penelitian}
\begin{itemize}
    \item \textbf{Judul Makalah:} \textit{GAN-Based Document Restoration with Frozen HTR Recognizer and Multi-Component Loss Optimization for Historical Manuscripts.}
    \item \textbf{Target Publikasi:} Jurnal Internasional - Advances in Human-Computer Interaction (Wiley Online Library).
    \item \textbf{Status Saat Ini:} \textbf{Submitted} (Tanggal Pengiriman: 17 November 2025).
\end{itemize}

% --- EVALUASI DIRI ---
\section*{E. Evaluasi Diri dan Rencana Tindak Lanjut}
\textbf{Capaian:}
\begin{itemize}
    \item Penelitian telah berhasil membuktikan hipotesis bahwa strategi \textit{Frozen Recognizer} lebih stabil dibandingkan \textit{Joint Training}, dengan penurunan CER signifikan pada dokumen uji.
    \item Draf tesis telah mencapai tahap finalisasi (V1) dengan kelengkapan Bab I hingga VI, termasuk analisis statistik hasil eksperimen.
    \item Makalah telah dikirimkan (submitted) ke jurnal target.
\end{itemize}

\textbf{Kendala \& Rencana:}
\begin{itemize}
    \item \textbf{Kendala:} Diperlukan ketelitian tinggi dalam pemformatan akhir dokumen (tabel, gambar, rumus) sesuai standar akademik ITB, serta persiapan materi presentasi yang ringkas namun komprehensif.
    \item \textbf{Rencana:} Melakukan revisi final berdasarkan masukan tanggal 28 November (terutama pada visualisasi alur metodologi), melaksanakan latihan presentasi seminar, dan memantau status \textit{review} jurnal.
\end{itemize}

% --- RUTINITAS ---
\section*{F. Rutinitas Mengerjakan Tesis}
\begin{itemize}
    \item \textbf{Pelatihan Model:} Menjalankan eksperimen pelatihan GAN menggunakan server GPU, memantau \textit{loss trajectory}, dan menyimpan \textit{checkpoint} model secara berkala.
    \item \textbf{Evaluasi Hasil:} Menganalisis metrik kualitas restorasi (PSNR, SSIM) dan metrik HTR (CER, WER) pada dataset uji untuk setiap konfigurasi eksperimen.
    \item \textbf{Studi Ablasi:} Melakukan eksperimen sistematis dengan menonaktifkan komponen tertentu (misal: \textit{perceptual loss}, \textit{CTC loss}) untuk memvalidasi kontribusi masing-masing.
    \item \textbf{Visualisasi \& Analisis:} Membuat visualisasi perbandingan hasil restorasi, \textit{confusion matrix} karakter, dan grafik konvergensi untuk interpretasi hasil.
    \item \textbf{Penulisan Dokumen:} Menyusun dan merevisi naskah tesis menggunakan LaTeX, termasuk pemformatan tabel, gambar, dan persamaan matematika.
    \item \textbf{Konsultasi Pembimbing:} Berdiskusi rutin dengan pembimbing untuk evaluasi progres, validasi metodologi, dan arahan perbaikan.
    \item \textbf{Manajemen Eksperimen:} Mencatat konfigurasi \textit{hyperparameter}, hasil eksperimen, dan temuan penting menggunakan sistem \textit{version control} (Git).
\end{itemize}

% --- HAMBATAN ---
\section*{G. Hambatan dalam Pengerjaan Tesis}
\begin{itemize}
    \item \textbf{Keterbatasan Komputasi:} Proses pelatihan model GAN memerlukan sumber daya GPU yang intensif, sehingga membutuhkan manajemen waktu eksperimen yang efisien.
    \item \textbf{Kompleksitas Integrasi:} Mengintegrasikan modul HTR (\textit{frozen recognizer}) dengan arsitektur GAN memerlukan pemahaman mendalam terhadap kedua domain.
    \item \textbf{Ketersediaan Data:} Dataset dokumen historis paleografi abad ke-16--18 yang berkualitas tinggi relatif terbatas dan memerlukan proses kurasi manual.
    \item \textbf{Pemformatan Akademik:} Penyesuaian format penulisan tesis sesuai standar ITB (tabel, gambar, rumus, sitasi) memerlukan waktu dan ketelitian ekstra.
    \item \textbf{Keseimbangan Waktu:} Menyeimbangkan waktu antara kegiatan penelitian, penulisan, dan aktivitas perkuliahan lainnya.
\end{itemize}

\vspace{1.5cm}

% --- TANDA TANGAN ---
\begin{table}[h!]
    \centering
    \begin{tabular}{p{7cm} p{7cm}}
        & Bandung, 12 Desember 2025 \\
        Menyetujui, & Mahasiswa, \\
        Pembimbing & \\
        \vspace{2.5cm} & \vspace{2.5cm} \\
        \textbf{Ir. I Gusti Bagus Baskara Nugraha, ST, MT, Ph.D} & \textbf{Jatniko Nur Mutaqin} \\
        NIP. ........................... & NIM. 23523314
    \end{tabular}
\end{table}

\end{document}